\documentclass[12pt,a4paper]{article}
\usepackage[utf8]{inputenc}
\usepackage[T1]{fontenc}
\usepackage[spanish]{babel}
\usepackage{amsmath,amssymb,amsfonts}
\usepackage{geometry}
\usepackage{booktabs}
\usepackage{hyperref}
\geometry{margin=2.5cm}

\title{Formulaci\'on Gauge-Invariante de la Perturbaci\'on del Horizonte y su Relaci\'on con la Perturbaci\'on de Curvatura Com\'ovil}
\author{Kevin Mu\~noz}
\date{}

\begin{document}
\maketitle

\begin{abstract}
La formalización de la dinámica cosmológica centrada en el horizonte identifica la perturbación del horizonte $\psi = \delta R_A / R_A$ con la perturbación de curvatura comóvil $\mathcal{R}$, pero deja la derivación gauge-invariante de esta identificación como un problema abierto. Este documento cierra esa brecha. Trabajando en gauge de Newton, calculamos $\psi$ a partir de la condición de horizonte aparente perturbado y realizamos la transformación de gauge al gauge comóvil, donde $\mathcal{R}$ está definida. La relación exacta en escalas super-horizonte es
\[
\psi = \frac{\varepsilon}{1 + \varepsilon}\,\mathcal{R}
\]
donde $\varepsilon = -\dot{H}/H^2$ es el parámetro de rodadura lenta. Para inflación de rodadura lenta ($\varepsilon \ll 1$), $\psi \approx \varepsilon\,\mathcal{R}$, de modo que los espectros de potencias primordiales satisfacen $P_\mathcal{R}(k) \approx P_\psi(k)/\varepsilon^2$. La forma espectral es idéntica para $\psi$ y $\mathcal{R}$. Este documento reemplaza la identificación esquemática $\mathcal{R} \sim \psi$ del Documento~5 por un resultado derivado y gauge-invariante, y corrige la ecuación (21) de ese documento.
\end{abstract}

\section{Motivación}

En la formalización (Documento~5), la perturbación de curvatura comóvil se introduce mediante la identificación esquemática
\begin{equation}
\mathcal{R} \sim \psi \tag{1}
\end{equation}
con la observación de que ``un tratamiento gauge-invariante completo sigue siendo un problema abierto.'' El espectro de potencias fue establecido como $P_\mathcal{R} = P_\psi$, lo cual solo es correcto salvo factores que involucran $\varepsilon$.

El presente documento deriva la relación gauge-invariante precisa entre $\psi$ y $\mathcal{R}$.

\section{Configuración: Perturbaciones Escalares en Gauge de Newton}

Consideramos un fondo FRW plano perturbado por modos escalares. En el \textbf{gauge de Newton (longitudinal)}, la métrica toma la forma
\begin{equation}
ds^2 = -(1 + 2\Phi)dt^2 + a^2(t)(1 - 2\Phi)\delta_{ij}\,dx^i dx^j \tag{2}
\end{equation}

donde $\Phi = \Psi$ para un fluido perfecto sin tensión anisotrópica. Las ecuaciones de fondo son
\begin{equation}
H^2 = \frac{8\pi G}{3}\rho, \qquad \dot{H} = -4\pi G(\rho + P) = -\varepsilon H^2 \tag{3}
\end{equation}

A primer orden, la perturbación métrica satisface la ligadura hamiltoniana
\begin{equation}
3H(\dot{\Phi} + H\Phi) + \frac{k^2}{a^2}\Phi = -4\pi G\,\delta\rho \tag{4}
\end{equation}

y la ligadura de momento
\begin{equation}
\dot{\Phi} + H\Phi = -4\pi G(\rho + P)\,V \tag{5}
\end{equation}

\section{La Perturbación del Horizonte en Gauge de Newton}

La tasa de expansión local efectiva es
\begin{equation}
H_\text{loc}(\mathbf{x}, t) = H + \dot{\Phi} - H\Phi + \frac{k^2 V}{3a} \tag{6}
\end{equation}

La perturbación fraccional del radio del horizonte aparente es
\begin{equation}
\psi \equiv \frac{\delta R_A}{R_A} = -\frac{\delta H}{H} = -\frac{\dot{\Phi} - H\Phi}{H} - \frac{k^2 V}{3aH} \tag{7}
\end{equation}

\textbf{En escalas super-horizonte} ($k \ll aH$), el término de velocidad está suprimido:
\begin{equation}
\psi \approx -\frac{\dot{\Phi}}{H} + \Phi \qquad (k \ll aH) \tag{8}
\end{equation}

Para el modo creciente, $\dot{\Phi} \to 0$, por lo que
\begin{equation}
\psi \approx \Phi \qquad \text{(super-horizonte, gauge de Newton)} \tag{9}
\end{equation}

\section{La Perturbación de Curvatura Comóvil}

La perturbación de curvatura comóvil se define como
\begin{equation}
\mathcal{R} = \Phi - \frac{H}{\dot{H}}\left(\dot{\Phi} + H\Phi\right) \tag{10}
\end{equation}

Expandiendo con $\dot{H} = -\varepsilon H^2$ y tomando el límite super-horizonte $\dot{\Phi} \to 0$:
\begin{equation}
\mathcal{R} \approx \Phi\,\frac{1 + \varepsilon}{\varepsilon} \tag{11}
\end{equation}

\section{Relación Gauge-Invariante Entre $\psi$ y $\mathcal{R}$}

Combinando las ecuaciones (9) y (11):
\begin{equation}
\psi \approx \Phi = \frac{\varepsilon}{1+\varepsilon}\,\mathcal{R} \tag{12}
\end{equation}

Reordenando:
\begin{equation}
\mathcal{R} = \frac{1 + \varepsilon}{\varepsilon}\,\psi \tag{13}
\end{equation}

Para hacer explícita la dependencia en $\dot{\Phi}$, sustituir $\Phi = \psi + \dot{\Phi}/H$ (de la ec.~(8)) en la ec.~(11):
\begin{equation}
\mathcal{R} = \frac{1+\varepsilon}{\varepsilon}\,\psi + \frac{\dot\Phi}{H} \tag{14}
\end{equation}
En el límite del modo creciente $\dot{\Phi} \to 0$, el último término se anula y se recupera la ec.~(13). $\square$

Este es el resultado central, válido en escalas super-horizonte para el modo creciente en el caso adiabático de fluido único.

\section{Casos Límite}

\paragraph{Inflación de rodadura lenta ($\varepsilon \ll 1$).}
\begin{equation}
\psi \approx \varepsilon\,\mathcal{R} \qquad \Longrightarrow \qquad \mathcal{R} \approx \frac{\psi}{\varepsilon} \tag{15}
\end{equation}

\paragraph{Dominación de radiación ($\varepsilon = 2$).}
\begin{equation}
\psi = \frac{2}{3}\,\mathcal{R} \tag{16}
\end{equation}

\paragraph{Dominación de materia ($\varepsilon = 3/2$).}
\begin{equation}
\psi = \frac{3}{5}\,\mathcal{R} \tag{17}
\end{equation}

\section{Derivación por Transformación de Gauge}

Bajo una reparametrización temporal $t \to t + \delta t(\mathbf{x}, t)$:
\begin{equation}
\Phi \to \Phi - H\,\delta t, \qquad \Psi_\text{metric} \to \Psi_\text{metric} + H\,\delta t \tag{18}
\end{equation}

Para llegar al gauge comóvil ($V_\text{com} = 0$), elegimos $\delta t = -V/a$. En este gauge:
\begin{equation}
\mathcal{R} = \Phi_\text{Newton} + H \cdot \frac{V}{a} \tag{19}
\end{equation}

De la ligadura de momento (5): $HV/a = -(\dot{\Phi} + H\Phi)/(\varepsilon H^2)$. Sustituyendo:
\begin{equation}
\mathcal{R} = \Phi\left(\frac{1+\varepsilon}{\varepsilon}\right) - \frac{\dot{\Phi}}{\varepsilon H} \tag{20}
\end{equation}

En escalas super-horizonte ($\dot{\Phi} \to 0$), con $\psi \approx \Phi$, esto recupera (13). $\square$

\section{Comportamiento Sub-Horizonte y Ecuación de Poisson}

En escalas sub-horizonte ($k \gg aH$):
\begin{equation}
\Phi \approx 3\left(\frac{aH}{k}\right)^2 \psi \to 0 \quad (k \gg aH) \tag{21}
\end{equation}

La identificación (12) se desintegra para $k \sim aH$. La relación relevante para la física del CMB es siempre el valor primordial establecido en la salida del horizonte.

\section{Implicación para los Espectros de Potencias}

De la relación gauge-invariante (13):
\begin{equation}
P_\mathcal{R}(k) = \left(\frac{1+\varepsilon}{\varepsilon}\right)^2 P_\psi(k) \tag{22}
\end{equation}

Para inflación de rodadura lenta ($\varepsilon \ll 1$):
\begin{equation}
P_\mathcal{R}(k) \approx \frac{P_\psi(k)}{\varepsilon^2}, \qquad \Delta_\mathcal{R}^2 = \frac{H^2}{2\pi^2 c_s \varepsilon^2} \tag{23}
\end{equation}

El \textbf{tilt espectral} no cambia:
\begin{equation}
n_s - 1 = \frac{d\ln P_\mathcal{R}}{d\ln k} = \frac{d\ln P_\psi}{d\ln k} \tag{24}
\end{equation}

\section{Correcciones al Documento 5}

\paragraph{Ecuación (18): $\mathcal{R} \sim \psi$.} Reemplazada por el resultado gauge-invariante exacto:
\begin{equation}
\mathcal{R} = \frac{1+\varepsilon}{\varepsilon}\,\psi \tag{25}
\end{equation}

\paragraph{Ecuación (21): $P_\mathcal{R}(k) = P_\psi(k)$.} Reemplazada por:
\begin{equation}
P_\mathcal{R}(k) = \left(\frac{1+\varepsilon}{\varepsilon}\right)^2 P_\psi(k) \approx \frac{P_\psi(k)}{\varepsilon^2} \quad (\varepsilon \ll 1) \tag{26}
\end{equation}

\section{Interpretación Física}

El factor $\varepsilon/(1+\varepsilon)$ en (13) tiene un significado físico directo en el marco centrado en el horizonte.

El parámetro de rodadura lenta $\varepsilon = -\dot{H}/H^2 = -\dot{R}_A/(R_A H)$ mide la tasa fraccional de cambio del radio del horizonte por tiempo de Hubble. Cuando $\varepsilon \ll 1$, el horizonte está casi congelado: una curvatura espacial dada $\mathcal{R}$ corresponde a solo una pequeña variación fraccional $\psi \approx \varepsilon \mathcal{R}$ en el radio del horizonte aparente.

Equivalentemente: el campo de horizonte $\phi \sim \psi$ es la amplitud de las fluctuaciones de borde, mientras que $\mathcal{R}$ es la amplitud de las fluctuaciones de curvatura en el bulk. El cociente $\mathcal{R}/\psi = (1+\varepsilon)/\varepsilon$ cuantifica cuánto amplifica el bulk la señal del borde. La amplificación es grande precisamente cuando el horizonte de fondo evoluciona lentamente --- consistente con la intuición de que un fondo casi-de Sitter actúa como un buen amplificador de fluctuaciones cuánticas del borde.

\section{Problema Abierto Restante}

La derivación está restringida al caso adiabático de fluido único y al límite super-horizonte. Extensiones a $k$ finito y a perturbaciones multi-fluido permanecen abiertas.

\section{Resumen}

\begin{table}[h]
\centering
\begin{tabular}{lll}
\toprule
Resultado & Expresión & Dominio de validez \\
\midrule
Relación gauge-invariante & $\psi = \frac{\varepsilon}{1+\varepsilon}\,\mathcal{R}$ & Super-horizonte, adiabático, fluido único \\
Límite de rodadura lenta & $\mathcal{R} \approx \psi/\varepsilon$ & $\varepsilon \ll 1$ \\
Relación de espectros & $P_\mathcal{R} = \frac{(1+\varepsilon)^2}{\varepsilon^2}\,P_\psi$ & Super-horizonte, primordial \\
Tilt espectral & $n_s(\mathcal{R}) - 1 = n_s(\psi) - 1 = 3 - 2\nu$ & Todas las escalas, orden líder \\
\bottomrule
\end{tabular}
\end{table}

\begin{thebibliography}{9}
\bibitem{bardeen1980} J. M. Bardeen, ``Gauge invariant cosmological perturbations,'' \textit{Phys. Rev. D} \textbf{22}, 1882 (1980).
\bibitem{kodama1984} H. Kodama y M. Sasaki, ``Cosmological perturbation theory,'' \textit{Prog. Theor. Phys. Suppl.} \textbf{78}, 1 (1984).
\bibitem{mukhanov1992} V. F. Mukhanov, H. A. Feldman y R. H. Brandenberger, ``Theory of cosmological perturbations,'' \textit{Phys. Rep.} \textbf{215}, 203 (1992).
\bibitem{weinberg2008} S. Weinberg, \textit{Cosmology}, Oxford University Press (2008).
\end{thebibliography}

\end{document}
